% Options for packages loaded elsewhere
\PassOptionsToPackage{unicode}{hyperref}
\PassOptionsToPackage{hyphens}{url}
\PassOptionsToPackage{space}{xeCJK}
\documentclass[
]{article}
\usepackage{xcolor}
\usepackage{amsmath,amssymb}
\setcounter{secnumdepth}{-\maxdimen} % remove section numbering
\usepackage{iftex}
\ifPDFTeX
  \usepackage[T1]{fontenc}
  \usepackage[utf8]{inputenc}
  \usepackage{textcomp} % provide euro and other symbols
\else % if luatex or xetex
  \usepackage{unicode-math} % this also loads fontspec
  \defaultfontfeatures{Scale=MatchLowercase}
  \defaultfontfeatures[\rmfamily]{Ligatures=TeX,Scale=1}
\fi
\usepackage{lmodern}
\ifPDFTeX\else
  % xetex/luatex font selection
    \setmainfont[Path=/usr/share/tor-browser/Browser/fonts/]{NotoSansJP-Regular.otf}
    \setsansfont[Path=/usr/share/tor-browser/Browser/fonts/]{NotoSansJP-Regular.otf}
    \setmonofont[Path=/usr/share/tor-browser/Browser/fonts/]{NotoSansJP-Regular.otf}
  \ifXeTeX
    \usepackage{xeCJK}
    \setCJKmainfont[Path=/usr/share/tor-browser/Browser/fonts/]{NotoSansJP-Regular.otf}
          \fi
  \ifLuaTeX
    \usepackage[]{luatexja-fontspec}
    \setmainjfont[Path=/usr/share/tor-browser/Browser/fonts/]{NotoSansJP-Regular.otf}
  \fi
\fi
% Use upquote if available, for straight quotes in verbatim environments
\IfFileExists{upquote.sty}{\usepackage{upquote}}{}
\IfFileExists{microtype.sty}{% use microtype if available
  \usepackage[]{microtype}
  \UseMicrotypeSet[protrusion]{basicmath} % disable protrusion for tt fonts
}{}
\makeatletter
\@ifundefined{KOMAClassName}{% if non-KOMA class
  \IfFileExists{parskip.sty}{%
    \usepackage{parskip}
  }{% else
    \setlength{\parindent}{0pt}
    \setlength{\parskip}{6pt plus 2pt minus 1pt}}
}{% if KOMA class
  \KOMAoptions{parskip=half}}
\makeatother
\usepackage{longtable,booktabs,array}
\usepackage{calc} % for calculating minipage widths
% Correct order of tables after \paragraph or \subparagraph
\usepackage{etoolbox}
\makeatletter
\patchcmd\longtable{\par}{\if@noskipsec\mbox{}\fi\par}{}{}
\makeatother
% Allow footnotes in longtable head/foot
\IfFileExists{footnotehyper.sty}{\usepackage{footnotehyper}}{\usepackage{footnote}}
\makesavenoteenv{longtable}
\usepackage{graphicx}
\makeatletter
\newsavebox\pandoc@box
\newcommand*\pandocbounded[1]{% scales image to fit in text height/width
  \sbox\pandoc@box{#1}%
  \Gscale@div\@tempa{\textheight}{\dimexpr\ht\pandoc@box+\dp\pandoc@box\relax}%
  \Gscale@div\@tempb{\linewidth}{\wd\pandoc@box}%
  \ifdim\@tempb\p@<\@tempa\p@\let\@tempa\@tempb\fi% select the smaller of both
  \ifdim\@tempa\p@<\p@\scalebox{\@tempa}{\usebox\pandoc@box}%
  \else\usebox{\pandoc@box}%
  \fi%
}
% Set default figure placement to htbp
\def\fps@figure{htbp}
\makeatother
\usepackage{svg}
\setlength{\emergencystretch}{3em} % prevent overfull lines
\providecommand{\tightlist}{%
  \setlength{\itemsep}{0pt}\setlength{\parskip}{0pt}}
\usepackage{bookmark}
\IfFileExists{xurl.sty}{\usepackage{xurl}}{} % add URL line breaks if available
\urlstyle{same}
\hypersetup{
  pdftitle={THP-CLEAR 政策提言書（最終版）},
  hidelinks,
  pdfcreator={LaTeX via pandoc}}

\title{THP-CLEAR 政策提言書（最終版）}
\author{}
\date{}

\begin{document}
\maketitle

\section{THP-CLEAR
政策提言書（統合版）}\label{thp-clear-ux653fux7b56ux63d0ux8a00ux66f8ux7d71ux5408ux7248}

\subsection{カバーレター}\label{ux30abux30d0ux30fcux30ecux30bfux30fc}

\textbf{件名：}
ステーブルコインを利用した決済インフラ整備に関する政策提言の送付について\\
\textbf{宛先：}
内閣官房（総理官邸）／自民党政務調査会／総務省／経済産業省／財務省 御中

拝啓\\
複合災厄下における我が国の決済継続性確保のため、\textbf{既存通貨に基づくステーブルコイン（フィアット完全準拠）を前提とする国家監査レイヤ}の導入を中核とした政策提言書をご送付申し上げます。\\
本書（統合版）は、政府主導で\textbf{THP-CLEAR}を全国配備し、\textbf{Google
/ Apple / 主要金融機関（SMBC
等）}と連携して\textbf{電子決済の即時性と監査可能性}を兼ね備えた決済網を構築するための実装計画です。

迅速な実証と制度整備をご検討賜りたく、何卒よろしくお願い申し上げます。

敬具

------\\
提出者：〇〇〇〇（所属・役職）\\
連絡先：xxxx@xxxx\\
日付：2025-10-13\\
------

\subsection{政策提言本文}\label{ux653fux7b56ux63d0ux8a00ux672cux6587}

\subsubsection{1.
提言の趣旨と背景}\label{ux63d0ux8a00ux306eux8da3ux65e8ux3068ux80ccux666f}

世界経済は金価格の急騰、米国内政の混乱、州債・市債市場の脆弱化など複合的要因により不安定化し、\textbf{米国決済網の停止を含む最悪シナリオ}が現実味を帯びている。\\
日本は依然としてドル建て決済への依存度が高く、ドル流通の停滞は即座に電力・物流・食料調達に影響するため、\textbf{円による即時決済が可能な代替ネットワーク}を早急に整備する必要がある。

\subsubsection{2. 提言の骨子}\label{ux63d0ux8a00ux306eux9aa8ux5b50}

本提言は、既に実装済みの分散監査型決済基盤 \textbf{THP-CLEAR}
を用い、\textbf{既存通貨に基づくステーブルコイン（フィアット完全準拠）を前提とする国家監査レイヤ}によって政府主導の即時決済網を整備する。

\begin{itemize}
\tightlist
\item
  \textbf{目的}

  \begin{itemize}
  \tightlist
  \item
    \textbf{既存通貨に基づくステーブルコイン（フィアット完全準拠）／E-20互換による即時決済網の確保}\\
  \item
    国家主導の\textbf{取引透明性・監査証跡（Witness）}の整備\\
  \item
    有事・停電時を含む\textbf{継続的な決済可用性}の維持（監査KPIはE-MAD連携の指標体系に準拠）
  \end{itemize}
\item
  \textbf{概要}

  \begin{itemize}
  \tightlist
  \item
    \textbf{JPYC（円建てステーブルコイン）}を軸に、SMBC、Google、Apple
    等と連携したバーチャルデビット決済を展開\\
  \item
    THP-CLEAR は決済情報の「When / Who / Whom / Coin /
    Value」の真実相当性のみを記録・監査し、各決済事業者の主張を第三者的に裏付ける\\
  \item
    AI（Ricoh
    等）が出力した判断を証跡で裏打ちし、\textbf{AI監査の唯一の担保}を提供する
  \end{itemize}
\end{itemize}

\subsubsection{3.
技術構成の概要}\label{ux6280ux8853ux69cbux6210ux306eux6982ux8981}

\begin{itemize}
\tightlist
\item
  \textbf{ハードウェア}

  \begin{itemize}
  \tightlist
  \item
    \textbf{8コア/16スレッド（TDP 35W 以下）}の省電力ミニ PC を採用\\
  \item
    OS は Linux ベース、ストレージは NVMe SSD（512GB～1TB）\\
  \item
    全国の自治体・市役所ネットワークに 1,000 ノードを分散配置\\
  \item
    設置は「置いて電源を入れるだけ」、異常時はリモート停止から交換
  \end{itemize}
\item
  \textbf{ソフトウェア}

  \begin{itemize}
  \tightlist
  \item
    \textbf{THP-CLEAR}：ステーブルコイン決済監査専用ソフトウェア（Go
    言語実装）\\
  \item
    暗号層に独自技術「ISE（恭順暗号）」を採用し、復号を不要化\\
  \item
    KMS（鍵管理サービス）は \textbf{独立した KMS
    Ring（3台以上／quorum=2）の独立クラスタ}で提供し、MQ
    とは別系統で運用する。Challenge
    段のみフェイルオーバーを許容し、Submit は同一 KMS
    エンドポイントに固定（セッション TTL／上限制御）
  \end{itemize}
\item
  \textbf{監査・運用}

  \begin{itemize}
  \tightlist
  \item
    監査情報は各ノードで即時ハッシュ化・署名され、東西 DC
    による二拠点クラスタで整合\\
  \item
    デジタル庁監督下でノード監視を行い、障害発生時は交換前提で迅速復旧\\
  \item
    すべての決済は\textbf{監査証跡を伴う即時性}を保持する
  \end{itemize}
\end{itemize}

\subsubsection{4. コスト試算（1,000
ノード構成）}\label{ux30b3ux30b9ux30c8ux8a66ux7b971000-ux30ceux30fcux30c9ux69cbux6210}

\begin{itemize}
\tightlist
\item
  MQ：8GB / NVMe 512GB\\
\item
  WN：4GB / NVMe 256GB（または下位）\\
\item
  DB：4GB / SSD 4TB×2 RAID1 TLC
\end{itemize}

\begin{longtable}[]{@{}llll@{}}
\toprule\noalign{}
構成 & 数量 & 単価 & 小計 \\
\midrule\noalign{}
\endhead
\bottomrule\noalign{}
\endlastfoot
MQ & 400 & 68,000 & 27,200,000 \\
WN & 500 & 62,000 & 31,000,000 \\
DB & 100 & 135,000 & 13,500,000 \\
バッテリ付 AC & 1,000 & 6,000 & 6,000,000 \\
3年 NBD 保守 & 1,000 & 10,000 & 10,000,000 \\
\textbf{合計（税抜）} & & & \textbf{87,700,000} \\
\textbf{消費税（10\%）} & & & 8,770,000 \\
\textbf{総計（税込）} & & & \textbf{96,470,000（約0.965億円）} \\
\end{longtable}

\begin{quote}
メインフレーム保守料（年間）約 10〜20 億円と比較し、1/10
以下のコストで同等以上の信頼性と可用性を確保。
\end{quote}

\subsection{技術仕様（要点抜粋）}\label{ux6280ux8853ux4ed5ux69d8ux8981ux70b9ux629cux7c8b}

\subsubsection{1. 前提}\label{ux524dux63d0}

\begin{itemize}
\tightlist
\item
  \textbf{既存通貨に基づくステーブルコイン（フィアット完全準拠）を前提とする国家監査レイヤ}（例：JPYC）\\
\item
  CLEAR は決済実行を担わず、「When / Who / Whom / Coin /
  Value」の真実相当性を保証する
\end{itemize}

\subsubsection{2. E-20
互換性（要点）}\label{e-20-ux4e92ux63dbux6027ux8981ux70b9}

\begin{itemize}
\tightlist
\item
  E-20
  はステーブルコインの発行体クレジットと監査可能性を前提とする互換ガイドライン（本提言内定義）\\
\item
  CLEAR は第三者監査レイヤとして各決済事業者の主張を暗号証跡で裏付ける
\end{itemize}

\subsubsection{3. 暗号・KMS}\label{ux6697ux53f7kms}

\begin{itemize}
\tightlist
\item
  Ed25519 署名、TLS/mTLS\\
\item
  KMS（鍵管理サービス）は \textbf{独立した KMS
  Ring（3台以上／quorum=2）の独立クラスタ}として動作し、MQ
  とは別プロセスで提供。Challenge 段のみフェイルオーバーを許容し、Submit
  は同一 KMS エンドポイントに固定（セッション TTL／上限制御）\\
\item
  監査ログ指標：\texttt{session\_id\ /\ kms\_endpoint\ /\ challenge\_latency\_ms\ /\ submit\_latency\_ms\ /\ challenge\_failover\_count}
\end{itemize}

\subsubsection{4.
データレイヤ}\label{ux30c7ux30fcux30bfux30ecux30a4ux30e4}

\begin{itemize}
\tightlist
\item
  Witness（JSONL+Merkle）→ Materialized（SQLite）\\
\item
  DB は \textbf{4TB×2 RAID1（TLC）}、順次書込中心、OP 確保と fstrim 運用
\end{itemize}

\subsubsection{5. 運用}\label{ux904bux7528}

\begin{itemize}
\tightlist
\item
  監視：Prometheus / Grafana（\texttt{kms\_*} 指標）\\
\item
  交換前提、UPS 不要、電源はバッテリ付 AC\\
\item
  自治体ネット＋東西 DC のハイブリッド分散構成
\end{itemize}

\subsubsection{6.
基準ノード仕様（固定）}\label{ux57faux6e96ux30ceux30fcux30c9ux4ed5ux69d8ux56faux5b9a}

\begin{itemize}
\tightlist
\item
  MQ：8GB / NVMe 512GB\\
\item
  WN：4GB / NVMe 256GB（または下位）\\
\item
  DB：4GB / SSD 4TB×2 RAID1 TLC
\end{itemize}

\subsection{全体構成図（SVG画像）}\label{ux5168ux4f53ux69cbux6210ux56f3svgux753bux50cf}

\pandocbounded{\includesvg[keepaspectratio]{./assets/diagrams/thp-clear-architecture.svg}}\\
図1. 主要コンポーネントとデータフロー（KMS は独立 Ring）

\subsection{スケールアウト要約}\label{ux30b9ux30b1ux30fcux30ebux30a2ux30a6ux30c8ux8981ux7d04}

\subsubsection{目的}\label{ux76eeux7684}

\begin{itemize}
\tightlist
\item
  国家規模（全世界分散）の\textbf{連続稼働監査}を 1,000
  ノード級で成立させる設計検証\\
\item
  \textbf{KMS 多重化（quorum=2）}と \textbf{Challenge-bound
  セッション固定}の実運用性評価
\end{itemize}

\subsubsection{主要結果}\label{ux4e3bux8981ux7d50ux679c}

\begin{itemize}
\tightlist
\item
  \textbf{KMS クォーラム成立}：Challenge 段のみフェイルオーバー、Submit
  は同一 KMS 固定で安定\\
\item
  \textbf{可観測性}：\texttt{kms\_challenge\_failover\_total\ /\ kms\_submit\_retry\_total\ /\ kms\_sessions\_in\_use}
  などの指標を整備\\
\item
  \textbf{冗長性}：ノードは交換前提（UPS 不要）。自治体ネット分散＋東西
  DC で国家級の停止条件に引き上げ\\
\item
  \textbf{性能方針}：Raft は DC 内低遅延、自治体ノードは非同期
  append。P95 500ms 目標
\end{itemize}

\subsubsection{推奨構成（抜粋）}\label{ux63a8ux5968ux69cbux6210ux629cux7c8b}

\begin{itemize}
\tightlist
\item
  MQ：8GB / NVMe 512GB\\
\item
  WN：4GB / NVMe 256GB（または下位）\\
\item
  DB：4GB / SSD 4TB×2 RAID1 TLC
\end{itemize}

\subsubsection{運用ドリル}\label{ux904bux7528ux30c9ux30eaux30eb}

\begin{itemize}
\tightlist
\item
  KMS 1 台停止 → failover 増分、quorum 継続\\
\item
  KMS 2 台停止 → quorum failed、監査記録と警報で早期復旧
\end{itemize}

\subsection{標準予算（1,000ノード）}\label{ux6a19ux6e96ux4e88ux7b971000ux30ceux30fcux30c9}

\subsubsection{前提}\label{ux524dux63d0-1}

\begin{itemize}
\tightlist
\item
  \textbf{既存通貨に基づくステーブルコイン（フィアット完全準拠）を前提とする国家監査レイヤ}（例：JPYC）\\
\item
  基準機：HP Elite Mini 805 G8（Ryzen 7 5700GE / 8C16T / TDP 35W / OS
  レス。同等品で代替可）\\
\item
  構成：MQ 400 / WN 500 / DB 100。バッテリ付 AC は PSE
  認証済み官公庁実績品\\
\item
  MQ：8GB / NVMe 512GB\\
\item
  WN：4GB / NVMe 256GB（または下位）\\
\item
  DB：4GB / SSD 4TB×2 RAID1 TLC
\end{itemize}

\subsubsection{単価（税抜・目安）}\label{ux5358ux4fa1ux7a0eux629cux76eeux5b89}

\begin{longtable}[]{@{}lll@{}}
\toprule\noalign{}
区分 & 構成 & 単価 \\
\midrule\noalign{}
\endhead
\bottomrule\noalign{}
\endlastfoot
MQ & 8GB / NVMe 512GB & ¥68,000 \\
WN & 4GB / NVMe 256GB & ¥62,000 \\
DB & 4GB / SSD 4TB×2（RAID1, TLC） & ¥135,000 \\
バッテリ付 AC & PSE 認証・互換品 & ¥6,000 \\
3 年 NBD 保守 & On-site & ¥10,000 \\
\end{longtable}

\subsubsection{合計（税抜→税込）}\label{ux5408ux8a08ux7a0eux629cux7a0eux8fbc}

\begin{longtable}[]{@{}llll@{}}
\toprule\noalign{}
構成 & 数量 & 単価 & 小計 \\
\midrule\noalign{}
\endhead
\bottomrule\noalign{}
\endlastfoot
MQ & 400 & 68,000 & 27,200,000 \\
WN & 500 & 62,000 & 31,000,000 \\
DB & 100 & 135,000 & 13,500,000 \\
バッテリ付 AC & 1,000 & 6,000 & 6,000,000 \\
3 年 NBD 保守 & 1,000 & 10,000 & 10,000,000 \\
\textbf{合計（税抜）} & & & \textbf{87,700,000} \\
\textbf{消費税（10\%）} & & & 8,770,000 \\
\textbf{総計（税込）} & & & \textbf{96,470,000（約0.965億円）} \\
\end{longtable}

\textbf{但し書き}\\
同等性能と保守体制を有する機器（DELL / ASUS 等）で代替可能。OS
レス（Linux）構成であり、Windows
ライセンス費は含まない。価格は時期・数量・為替によって変動し、競争入札での圧縮余地がある。

\subsection{人道特化モデル（100ノード）}\label{ux4ebaux9053ux7279ux5316ux30e2ux30c7ux30eb100ux30ceux30fcux30c9}

本モデルは THP-CLEAR
を中核とした\textbf{パキスタン--アフガニスタン救援向け人道支援決済インフラ}であり、国内安全拠点クラスタ
100 ノードで稼働させ、衛星と可搬 Wi-Fi により現地へ接続する。

\subsubsection{構成概要}\label{ux69cbux6210ux6982ux8981}

\begin{itemize}
\tightlist
\item
  配置：国内 DC（東西 2 拠点）に MQ40 / WN50 / DB10 の 100
  ノードを設置\\
\item
  通信：NEC / NTT の可搬 Wi-Fi 中継＋ポータブル電源、Starlink 衛星回線\\
\item
  決済：JPYC＋SMBC（口座・信用保証）、Google / Apple / Amazon Pay
  と連携、必要に応じ Tether 参入を要請\\
\item
  監査：THP-CLEAR Witness で「When / Who / Whom / Coin /
  Value」を記録し、UNHCR / WHO / NGO 連合へ開放\\
\item
  暗号：ISE（恭順暗号）＋ Ed25519 署名。KMS は \textbf{独立した KMS
  Ring（3台以上／quorum=2）の独立クラスタ}として国内 DC
  に常設し、Challenge 段のみフェイルオーバー、Submit は同一 KMS
  固定（セッション TTL／上限制御）\\
\item
  基準ノード仕様（固定）：

  \begin{itemize}
  \tightlist
  \item
    MQ：8GB / NVMe 512GB\\
  \item
    WN：4GB / NVMe 256GB（または下位）\\
  \item
    DB：4GB / SSD 4TB×2 RAID1 TLC
  \end{itemize}
\end{itemize}

\subsubsection{体制・役割}\label{ux4f53ux5236ux5f79ux5272}

\begin{itemize}
\tightlist
\item
  通信・電源：NEC / NTT\\
\item
  衛星通信：SpaceX（Starlink）\\
\item
  決済統合：Google / Apple / Amazon / SMBC / JPYC\\
\item
  監査基盤：THP-CLEAR Core（日本）\\
\item
  人道統制：UNHCR / WHO / NGO 連合\\
\item
  技術監修：THP-CLEAR 開発局（想定：デジタル庁内）
\end{itemize}

\subsubsection{運用目的と効果}\label{ux904bux7528ux76eeux7684ux3068ux52b9ux679c}

\begin{itemize}
\tightlist
\item
  \textbf{人道目的}：食料・医薬・教育・給与支援の取引と物流を即時かつ透明に管理\\
\item
  \textbf{技術目的}：現地ノード不要のリモート監査ネットワークを確立\\
\item
  \textbf{政策目的}：円建てステーブルコインによる国際人道決済標準化と AI
  監査（Ricoh 等）による不正抑止
\end{itemize}

\subsection{人道特化予算（100ノード）}\label{ux4ebaux9053ux7279ux5316ux4e88ux7b97100ux30ceux30fcux30c9}

\begin{itemize}
\tightlist
\item
  MQ：8GB / NVMe 512GB\\
\item
  WN：4GB / NVMe 256GB（または下位）\\
\item
  DB：4GB / SSD 4TB×2 RAID1 TLC
\end{itemize}

\begin{longtable}[]{@{}lrrr@{}}
\toprule\noalign{}
項目 & 数量 & 単価 & 小計 \\
\midrule\noalign{}
\endhead
\bottomrule\noalign{}
\endlastfoot
MQ（8C16T / 8GB / NVMe 512GB） & 40 & ¥68,000 & ¥2,720,000 \\
WN（8C16T / 4GB / NVMe 256GB） & 50 & ¥62,000 & ¥3,100,000 \\
DB（8C16T / 4GB / SSD 4TB×2 RAID1） & 10 & ¥135,000 & ¥1,350,000 \\
バッテリ付 AC（PSE 認証） & 100 & ¥6,000 & ¥600,000 \\
可搬 Wi-Fi 中継（NEC / NTT） & 一式 & --- & ¥3,000,000 \\
Starlink（端末×5＋回線） & 一式 & --- & ¥2,000,000 \\
Pixel 型落ち端末（Google 協力） & 200 & ¥10,000 & ¥2,000,000 \\
技術者教育・運用立上げ & 一式 & --- & ¥1,500,000 \\
初年度運用費（通信・監査・保守） & 一式 & --- & ¥2,000,000 \\
\textbf{合計（税抜）} & & & \textbf{¥18,270,000} \\
\textbf{消費税（10\%）} & & & \textbf{¥1,827,000} \\
\textbf{総計（税込）} & & & \textbf{¥20,097,000（約2,000万円）} \\
\end{longtable}

\subsubsection{体制・運用イメージ}\label{ux4f53ux5236ux904bux7528ux30a4ux30e1ux30fcux30b8}

\begin{itemize}
\tightlist
\item
  NEC / NTT：中継・電源\\
\item
  SpaceX：衛星回線\\
\item
  SMBC・JPYC：決済・口座\\
\item
  Google・Apple・Amazon：即時決済連携\\
\item
  UNHCR・WHO・NGO：現地統制\\
\item
  THP-CLEAR 開発局：監査・教育
\end{itemize}

\subsubsection{導入スケジュール（概略）}\label{ux5c0eux5165ux30b9ux30b1ux30b8ux30e5ux30fcux30ebux6982ux7565}

\begin{longtable}[]{@{}ll@{}}
\toprule\noalign{}
期間 & 内容 \\
\midrule\noalign{}
\endhead
\bottomrule\noalign{}
\endlastfoot
T0（決定日） & 政府・企業・NGO 合意 \\
T+30 日 & 調達・ライセンス・通信契約 \\
T+60 日 & 衛星／Wi-Fi 接続試験・JPYC 統合 \\
T+90 日 & 稼働開始・UN/NGO 立会監査開始 \\
\end{longtable}

\subsubsection{政策意義}\label{ux653fux7b56ux610fux7fa9}

\textbf{既存通貨に基づくステーブルコイン（フィアット完全準拠）を前提とする国家監査レイヤ}により、100
台・約 2,000
万円で安全・即時・透明な人道決済を実現。円建て人道支援モデルの国際標準化を先導し、デジタル円の社会的受容性を実証する。

\subsection{結語}\label{ux7d50ux8a9e}

THP-CLEAR は約 1
億円規模で導入できる\textbf{唯一の国家級ステーブルコイン監査アーキテクチャ}であり、ソフトウェアは既に実装済みで速やかな稼働が可能である。\\
本書（統合版）に示した国内標準配備と人道特化モデルを同時に進めることで、日本は\textbf{ポスト・フィアット時代の金融安定装置}を世界に先駆けて確立できる。\\
関係省庁の連携のもと早急な実証と本格運用をご検討いただきたい。

\subsection{付記（提出形態）}\label{ux4ed8ux8a18ux63d0ux51faux5f62ux614b}

本書は全章統合済みの最終版（単一 PDF）です。別紙は存在しません。

\end{document}
